% Flavor catalog per-process draft.
% process_id: L005
% family: charged_lepton
% owner_agent_id: PKA-L005
% writer_agent_id: null
% checker_agent_id: null
% opus_signoff_id: null
% Canonical metadata sidecar: L005.yaml
% Status is controlled by the YAML sidecar status_history, not this comment.

\section{L005: \texorpdfstring{\(\mu-e\)}{mu-e} Conversion in Titanium}

\subsection*{Process}
\textbf{Standard notation:} \(R_{\mu e}^{\rm Ti} =
\Gamma(\mu^-{\rm Ti}\to e^-{\rm Ti})/\Gamma_{\rm capture}(\mu^-{\rm Ti})\).
This entry covers coherent neutrinoless muon-to-electron conversion in a
titanium target, the SINDRUM II bound that sits between the aluminum program
entry and the stronger gold-target benchmark.

\subsection*{PDG / equivalent value(s)}
The streamed PDG live muon listing, generated on 2026-05-14 and accessed on
2026-05-16, gives
\[
  R_{\mu e}^{\rm Ti}<4.3\times10^{-12}\qquad (90\%\ {\rm C.L.}) .
\]
The PDG line identifies the source as DOHMEN 93, technique SPEC, comment
SINDRUM II, and notes that the quoted limit assumes conversion to the nuclear
ground state.  The local minimal snapshot is
\texttt{flavor\_catalog/references/L005/pdg2026\_muon\_listing\_s004\_ti\_conversion.txt};
the SINDRUM II publication metadata is snapshotted separately in
\texttt{flavor\_catalog/references/L005/sindrumii\_1993\_titanium\_crossref.txt}.

\subsection*{Relevance to the RS / anarchic-flavor pipeline}
In a warped lepton-flavor extension, \(\mu-e\) conversion probes lepton-quark
contact operators, flavor-changing \(Z\)-like couplings, scalar currents, and
dipoles.  It is therefore complementary to \(\mu\to e\gamma\), which tests only
the dipole lane in the current repository.  Titanium is useful historically
because it is a real published conversion target, but its interpretation is
not target-universal.

\subsection*{Post-2008 developments}
No newer titanium-target measurement was found in the PDG stream used here
after the arXiv:0804.1954 RS-flavor baseline.  The main development since that
baseline is experimental: Mu2e describes a coherent \(\mu^-N\to e^-N\)
program aiming roughly four orders of magnitude beyond the previous published
conversion limits, and COMET Phase-I quotes an aluminum sensitivity goal
\(3.1\times10^{-15}\) with an expected \(90\%\) upper limit
\(7.0\times10^{-15}\).  These are not titanium measurements; they explain why
the Ti number is now primarily a legacy benchmark and target-comparison input.

\subsection*{Constraint validity and model dependence}
The experimental statement is robust as a null search normalized to the
ordinary muon-capture rate in titanium.  The mapping to RS parameters is
model-dependent: dipole, vector, and scalar amplitudes can interfere, and the
nuclear overlap integrals and capture normalization are target-specific.
Direct comparison with gold or aluminum therefore requires an operator basis
and nuclear-convention choice.

\subsection*{Code coverage in this repo}
\textbf{NO.}  The required greps over \texttt{quarkConstraints/}, \texttt{qcd/},
\texttt{flavorConstraints/}, \texttt{neutrinos/}, \texttt{yukawa/},
\texttt{warpConfig/}, \texttt{solvers/}, \texttt{scanParams/}, and
\texttt{tests/} found no SINDRUM, Mu2e, COMET, \(R_{\mu e}\), titanium
conversion, or general \(\mu-e\)-conversion implementation.  The only adjacent
charged-LFV code is the dipole-only \(\mu\to e\gamma\) checker at
\texttt{flavorConstraints/muToEGamma.py}:75, called by the scan at
\texttt{scanParams/scan.py}:524.

\subsection*{Implementation difficulty}
\textbf{HIGH.}  A production constraint needs a new coherent conversion
observable, lepton-quark Wilson coefficients, titanium nuclear overlap and
capture inputs, and likely EFT matching/running shared with \(\mu\to e\gamma\)
and \(\mu\to3e\).  This is a new target-dependent mode calculation, not a new
limit for the existing dipole checker.

\subsection*{Key references}
Process-local source keys before bibliography consolidation:
\texttt{PDG2026\_MuonTiConversion}, \texttt{SINDRUMII1993\_TitaniumMuE},
\texttt{Mu2eTDR2015}, \texttt{COMETPhaseITDR2020}, and \texttt{CFW2008}.
